\begin{figure*}
    \includegraphics[width=\linewidth]{Figures/CATALOG.pdf}
    \caption{Comparison of the colour-redshift distributions for \texttt{OpenUniverse2024} (orange contours) and \texttt{CosmoDC2} (green contours) at LSST Y10 depth, showing confidence contours at \(0.5\sigma, 1.0\sigma, 2.0\sigma,\) and \(3.0\sigma\). From top to bottom, the rows show the \(u - g\), \(g - r\), \(i - z\), and \(z - y\) colours, while the columns correspond to bins in \(i\)-band magnitude: \(16.0 < i < 21.0\), \(21.0 < i < 22.5\), \(22.5 < i < 24.0\), and \(24.0 < i < 25.3\). The contours highlight the relative coverage of colour-redshift space: \texttt{CosmoDC2} generally spans a broader region, while there remain areas occupied only by \texttt{OpenUniverse2024}, where resampling \texttt{CosmoDC2} galaxies alone cannot fully reconcile the discrepancy.}
    \label{fig:1}
\end{figure*}

To ensure realism and utility, tests of the data‐augmentation method using simulated galaxies must employ catalogues that exhibit authentic differences. In practice, the discrepancy between the real observation and one simulation should be comparable to that between two distinct simulations employed in our study. If the mocks were identical, augmentation would trivially succeed but offer no insight into its effectiveness in real analyses; if they were unrealistically different, augmentation would fail outright. 

In this section, we outline the key features of the two input mock galaxy catalogues, \texttt{CosmoDC2} and \texttt{OpenUniverse2024}, that serve as the starting points for our data manipulation and analysis. The \texttt{OpenUniverse2024} catalogue is utilised to emulate the properties of galaxies expected in forthcoming LSST observations, whereas \texttt{CosmoDC2} is used here solely for data augmentation. Our selected catalogues closely resemble actual observations; nevertheless, they exhibit differences, offering an appropriate testing environment for the intended application of the method. This variation effectively represents the systematic discrepancies between synthetic catalogues and real observations, thereby reproducing a significant source of bias in photometric-redshift estimation through data augmentation.

\subsection{\texttt{CosmoDC2}}

The \texttt{CosmoDC2} catalogue \citep{2019ApJS..245...26K} is a comprehensive, purpose‐designed resource for \textit{Rubin} cosmology. This constitutes the foundation of the second Data Challenge \citep[DC2;][]{2021ApJS..253...31L} conducted by the LSST DESC. Derived from the Outer Rim N-body simulation \citep{2019ApJS..245...16H}, it models a cosmic volume of \( \left( 4.225\, \mathrm{Gpc} \right)^3 \) using over a trillion \(10240^3\) particles, each with a mass of \(2.6 \times 10^9 \, M_{\odot}\). The simulation suite comprises roughly \(100\) temporal snapshots, recorded from \(z \sim 10\) to \(z = 0\), and was created using the Hybrid/Hardware Accelerated Cosmology Code (HACC; \citealp{2016NewA...42...49H}) on the IBM BG/Q system Mira at the Argonne Leadership Computing Facility (ALCF\footnote{\url{https://www.alcf.anl.gov}}).

To incorporate realistic galaxy properties, \texttt{CosmoDC2} employs \texttt{GalSampler} \citep{2020MNRAS.495.5040H}, a hybrid approach that merges semi-analytic modelling using \texttt{Galacticus} \citep{2012MNRAS.419.3590B}, with empirical techniques based on the \texttt{UniverseMachine} \citep{2019MNRAS.488.3143B}. This method transfers sophisticated stellar population synthesis prescriptions to cosmological simulations via halo-to-halo correspondence, incorporating more than \(500\) galaxy attributes, including the \textit{Rubin} photometry in the \(u\), \(g\), \(r\), \(i\), \(z\), \(y\) filters, morphology, halo properties, and weak lensing shear. \texttt{CosmoDC2} has been validated against the scientific objectives of LSST DESC \citep{2022OJAp....5E...1K}, and its simulated catalogues span over \(440 \,\mathrm{deg}^2\) of sky up to a redshift of \(z = 3\), reaching a magnitude depth of \(r = 28 \) and supporting a broad range of scientific applications.

\subsection{\texttt{OpenUniverse2024}}

The \texttt{OpenUniverse2024} extragalactic catalogue \citep{2025MNRAS.544.3799O} builds upon the same Outer Rim dark‐matter simulation used in \texttt{CosmoDC2}, but uses the differentiable \texttt{Diffsky} framework to populate galaxies within the same cosmological large‐scale structure, with details introduced in \href{https://lsstdesc-diffsky.readthedocs.io/en/latest/index.html}{\texttt{lsstdesc-diffsky}}. Rather than relying on discrete merger‐tree assembly histories and star‐formation tracks, \texttt{Diffsky} utilises the auto‐differentiable modules \texttt{DiffMAH} \citep{2021OJAp....4E...7H}, \texttt{DiffStar} \citep{2023MNRAS.518..562A}, and \texttt{DSPS} \citep{2023MNRAS.521.1741H} to derive smooth, physically motivated parametrisations of mass assembly histories, star formation histories, and stellar population synthesis modelling. This approach yields robust predictions of stellar mass and star formation rates, even in poorly resolved haloes. The resulting position, lensing shear and convergence, morphological parameters, and multi‐wavelength SED of each galaxy are consolidated within \texttt{SkyCatalog}, a unified truth‐catalogue format. \texttt{OpenUniverse2024} covers approximately \(110\,\mathrm{deg}^2\) -- corresponding to both the LSST Deep Drilling Fields (DDFs; \citealp{2024ApJS..275...21G}) and the Roman High-Latitude Imaging Survey \citep{2019arXiv190205569A}.

\subsection{Comparison of catalogues}

Both \texttt{CosmoDC2} and \texttt{OpenUniverse2024} employ galaxy SED models that have been rigorously validated against deep, multi-band survey data. Owing to its fully differentiable framework, the mock colours and number counts of \texttt{OpenUniverse2024} agree with those of the COSMOS2020 survey at the per cent level \citep{2022ApJS..258...11W, 2022A&A...664A..61S}. Similarly, \texttt{CosmoDC2} demonstrates excellent consistency in its bright, low-redshift galaxy population when evaluated with the LSST DESCQA validation suite \citep{2018ApJS..234...36M, 2019ApJS..245...26K, 2021ApJS..253...31L}.

As illustrated by the green contours in Figure~\ref{fig:1}, the colour–redshift distribution of \texttt{CosmoDC2} exhibits pronounced discreteness at high redshift compared to the orange contours for \texttt{OpenUniverse2024}. This behaviour originates from its dependence on a restricted array of stellar population templates. We therefore utilise \texttt{CosmoDC2} to construct our augmentation catalogues, thereby compensating for gaps in spectroscopic coverage. In a practical analysis, these augmented catalogues would likewise be simulated datasets, with their SEDs mildly differing from those in the observational datasets.

Conversely, the auto-differentiable, parameterised models utilised by \texttt{OpenUniverse2024} produce smoother and more physical colour distributions, rendering it an excellent proxy for LSST-like observations. In this study, we employ \texttt{OpenUniverse2024}, a high-fidelity LSST-like mock catalogue, as a proxy for the observational dataset expected from the forthcoming survey.

At low redshift (\(z_{\mathrm{true}} \lesssim 1.5\)), both catalogues exhibit excellent agreement in the \(i - z\) and \(z - y\) colours, particularly for bright galaxies (\(i \lesssim 22.5\)). Conversely, \texttt{CosmoDC2} shows a systematic excess of redder \(u - g\) and \(g - r\) colours relative to \texttt{OpenUniverse2024} especially among low-redshift (\(z_{\mathrm{true}} \lesssim 1.5\)) and faint galaxies (\(i \gtrsim 22.5\)), reflecting differences in the treatment of stellar populations at the blue end of the SED. As detailed below, these offsets can be partly mitigated by weighted resampling of the \texttt{CosmoDC2} galaxies to match the smoother distributions of \texttt{OpenUniverse2024}. 

At fainter magnitudes (\(i \gtrsim 22.5\)) and higher redshifts (\(z_{\mathrm{true}} \gtrsim 1.5\)), where spectroscopic calibrators are scarce, the two SED models diverge markedly, reducing their overlap. This divergence persists after resampling and thus represents a significant source of systematic error in photometric-redshift modelling. Figure~\ref{fig:1} illustrates these comparisons at LSST Y10 depth -- chosen to emphasise SED differences among the faintest galaxies (\(i\gtrsim24.0\)) -- although identical trends are evident in other observing scenarios.